%% AHAAR -- Kitchen System : IEEE-style Project Report
%% Compile with pdflatex, then bibtex, then pdflatex x2
%% Overleaf: upload main.tex, references.bib ; set compiler to pdfLaTeX
\documentclass[journal,onecolumn,12pt]{IEEEtran}

\usepackage[utf8]{inputenc}
\usepackage[T1]{fontenc}
\usepackage{graphicx}
\usepackage{amsmath,amssymb}
\usepackage{booktabs}
\usepackage{array}
\usepackage{url}
\usepackage{hyperref}
\usepackage{cite}
\usepackage{xcolor}
\usepackage{listings}
\usepackage{tabularx}

\hypersetup{colorlinks=true,linkcolor=blue,citecolor=blue,urlcolor=blue}

\lstset{
  basicstyle=\ttfamily\footnotesize,
  breaklines=true,
  columns=fullflexible,
  frame=single,
  showstringspaces=false
}

\begin{document}

\title{AHAAR: An AI-Driven Household Kitchen Assistant for Indian Cuisine with CGM-Aware Nutritional Personalization}

\author{%
  Rajiv Agarwal, Alan Joseph, Aryan Raina\\
  \textit{Group 19, Academic Year 2026--27}\\
  \textit{Project Guide: Ms.\ Sharlene Rebeiro}%
  \thanks{Manuscript prepared for the AHAAR Kitchen System project. Source code: \url{https://github.com/workspace220106/Aahar-Kitchen-System}}
}

\markboth{AHAAR Project Report, 2026--27}{Agarwal \MakeLowercase{\textit{et al.}}: AHAAR -- Kitchen System}

\maketitle

%====================================================================
\begin{abstract}
Household food waste, repetitive cooking decisions, and the lack of
personalized dietary guidance remain persistent problems in Indian
kitchens. We present \textbf{AHAAR}, a multi-page web application that
turns a single photograph of the inside of a refrigerator into a ranked
set of culturally-appropriate recipes, portioned per family member and
re-ranked over time by each user's personal glycemic response captured
through a Continuous Glucose Monitor (CGM). The frontend is already
implemented in React, Vite, and TailwindCSS; this report proposes the
complete backend, machine-learning, and database stack needed to
operationalize it. The core technical contributions are: (1) a
fridge-interior Indian ingredient detection dataset and a YOLOv11-based
detector trained on it; (2) a cross-modal recipe-embedding index built
over scraped Indian recipes, replacing hand-coded ingredient lists;
(3) a hybrid nutrition-aware recipe recommender whose weights adapt to
each family member's health profile; (4) a reference-object portion
estimator that avoids the hardware cost of RGB-D sensors; and (5) a
per-user CGM-aware recipe re-ranker inspired by the Weizmann Institute's
personalized nutrition work. We ground every design choice in ten
peer-reviewed IEEE papers spanning object detection, cross-modal recipe
retrieval, nutrition-aware recommendation, portion estimation, and
glucose-driven dietary monitoring. We describe the Java~26 Spring Boot
backend, PostgreSQL schema, and Deep Java Library (DJL) training
pipeline that realise these contributions, and present the dataset
construction, evaluation metrics, and ablation plan that will produce
the final results.
\end{abstract}

\begin{IEEEkeywords}
Food recognition, YOLO, Indian cuisine, recipe recommendation,
continuous glucose monitoring, personalized nutrition, deep learning,
Spring Boot, Deep Java Library, dietary assessment.
\end{IEEEkeywords}

%====================================================================
\tableofcontents
\vspace{1em}
\hrule
\vspace{1em}

%====================================================================
\section{Introduction}
\IEEEPARstart{T}{he} question \textit{``what's for dinner?''} is asked
in roughly 300~million Indian households every day. The answer is
constrained by what is physically in the refrigerator, by cultural and
regional cuisine preferences, by the age and health status of each
family member, and -- for the roughly 101~million Indians living with
diabetes -- by the glycemic impact of the chosen meal. No existing
application addresses all four constraints simultaneously for Indian
cuisine.

AHAAR is an attempt to do exactly this. It is organized as a sequence
of pages that mirror the real decision flow in a kitchen:
\textit{Fridge Scan $\rightarrow$ Detection $\rightarrow$ Pantry $\rightarrow$ Recipes $\rightarrow$ Family Portions $\rightarrow$ Cook-along $\rightarrow$ Health}.
The frontend already implements this flow; what remains is to replace
the demonstration data with real machine-learning inference, a
persistent database, a trained ingredient detector, a recipe retrieval
index, a nutrition-aware ranker, and -- as the project's most
differentiated contribution -- a glucose-response-aware personalisation
layer.

This report consolidates: (i) the current system's architecture;
(ii) a critical review of ten IEEE papers that underpin our proposed
upgrades; (iii) the complete proposed methodology, including a
Java-26-based backend and ML training pipeline using the Deep Java
Library; (iv) the data-collection and evaluation protocol; and
(v) a discussion of the novelty, limitations, and future work.

\textbf{Contributions.} We claim four contributions:
\begin{enumerate}
  \item A \emph{fridge-interior} Indian ingredient dataset -- distinct
        from existing \emph{plated-meal} Indian food datasets
        \cite{foodnet2017,indianfood10,indianfoodnet2023}.
  \item A light-and-sharpness pre-filter that reduces label noise
        during data collection and inference failure during use.
  \item A unified multi-head architecture for detection,
        portion-weight regression, and nutrition lookup -- in the
        spirit of \cite{yolov11indian2025} but extended with recipe
        retrieval.
  \item A per-user glucose-response-aware recipe ranker for Indian
        cuisine, bridging the CGM-macronutrient literature
        \cite{huo2019cgm} with cross-modal recipe retrieval
        \cite{marin2021recipe1m}.
\end{enumerate}

%====================================================================
\section{Related Work}
\label{sec:related}

We organise the related work along the five technical sub-problems
addressed by AHAAR.

\subsection{Object Detection and Food Recognition}

\textbf{YOLO}~\cite{redmon2016yolo} reframed object detection as a
single regression problem, predicting bounding boxes and class
probabilities in one network pass. This real-time, end-to-end approach
is directly suitable for on-device fridge scanning, where latency
budgets rule out two-stage detectors.

\textbf{FoodNet}~\cite{foodnet2017}, published in \textit{IEEE Signal
Processing Letters}, proposed an ensemble of deep CNNs for food
classification and -- importantly for us -- introduced the first
Indian food image database (50 classes, 100 images per class). It
established the baseline that domain-specific ensembles beat generic
backbones on Indian cuisine.

\textbf{IndianFood10 / YOLOv4 Platters}~\cite{indianfood10} applied
transfer learning with YOLOv4 to a 10-class dataset of staple Indian
plates and reported 91.8\% mAP, framing the central difficulties as
``lack of annotated Indian food datasets, non-distinct boundaries
between dishes on one plate, and high intra-class variation.''

\textbf{YOLO V11m Indian Cuisine}~\cite{yolov11indian2025}, the most
recent IEEE paper in this space, combined 17-class Indian food
detection with portion-based nutritional estimation and reached
\(\text{mAP}@50 = 0.94\). Its two-stage pipeline (detect $\rightarrow$
look up nutrition) is the architectural template we extend.

\subsection{Cross-Modal Recipe Retrieval}

\textbf{im2recipe}~\cite{salvador2017im2recipe}, presented at CVPR
2017, learns a joint embedding of recipe text and food images so that
either modality can retrieve the other. This is the key idea that
replaces our current string-level matching
(\verb|uses: ['spinach']|) with semantically-robust retrieval
(\textit{spinach} $\approx$ \textit{palak} $\approx$ \textit{saag}).

\textbf{Recipe1M+}~\cite{marin2021recipe1m}, the IEEE TPAMI extension
of im2recipe, scaled this joint embedding to more than one million
recipes and 13 million images and reported near-human retrieval
performance. We adopt its training objective and dataset conventions.

\textbf{Inverse Cooking}~\cite{salvador2019inverse} (CVPR 2019)
predicts an unordered set of ingredients from an image and then
generates full cooking instructions. The set-prediction formulation
sidesteps per-item occlusion in a crowded fridge and inspires our
``Lazy mode'' feature.

\subsection{Nutrition-Aware Recommendation}

\textbf{Yera-Toledo \emph{et al.}}~\cite{yeratoledo2019access}, in
\textit{IEEE Access} 2019, proposed a hybrid food recommender that
combines content-based filtering with explicit nutritional constraints.
Its weighted scoring formulation underlies our per-family-member
re-ranking weights.

\subsection{Portion and Volume Estimation}

\textbf{Gao, Lo, and Lo}~\cite{gao2018bsn} (IEEE BSN 2018) estimated
food volume from a wearable camera using SLAM, a modified convex-hull
algorithm, and a 3D mesh reconstruction step. The paper is instructive
primarily for what it concludes: reliable monocular volume estimation
is \emph{hard} without either depth sensors or a visible reference
object. We adopt the latter.

\subsection{CGM-Driven Dietary Monitoring}

\textbf{Huo \emph{et al.}}~\cite{huo2019cgm} (IEEE BHI 2019) showed
that meal \emph{macronutrient composition} can be inferred from the
glucose curve that follows the meal. Together with the Zeevi
\emph{et al.} personalised-nutrition work from \textit{Cell}
(cited for context only), this is the scientific basis for our
\texttt{/health} page: GI is not universal, so AHAAR should rank
recipes by \emph{this user's} predicted glucose response, not a
textbook GI table.

\medskip
Table~\ref{tab:papermap} summarises how each paper maps to a
component of the AHAAR system.

\begin{table}[!t]
\caption{Mapping of IEEE references to AHAAR system components.}
\label{tab:papermap}
\centering
\footnotesize
\begin{tabularx}{\linewidth}{l l X}
\toprule
\textbf{Ref.} & \textbf{Venue} & \textbf{AHAAR component informed} \\
\midrule
\cite{redmon2016yolo}        & CVPR 2016        & \texttt{/scan}, \texttt{/detect} base detector \\
\cite{foodnet2017}           & IEEE SPL 2017    & Indian-food backbone, ensembling \\
\cite{indianfood10}          & IEEE ICDEW 2022  & Pre-training dataset, YOLO transfer \\
\cite{yolov11indian2025}     & IEEE CICN 2025   & Detection + nutrition pipeline \\
\cite{salvador2017im2recipe} & IEEE CVPR 2017   & Recipe embedding (\texttt{/recipes}) \\
\cite{marin2021recipe1m}     & IEEE TPAMI 2021  & Recipe1M+ scale, loss \\
\cite{salvador2019inverse}   & IEEE CVPR 2019   & ``Lazy mode'', set prediction \\
\cite{yeratoledo2019access}  & IEEE Access 2019 & Nutrition-weighted ranking \\
\cite{gao2018bsn}            & IEEE BSN 2018    & Portion sizing design choice \\
\cite{huo2019cgm}            & IEEE BHI 2019    & \texttt{/health} CGM re-ranker \\
\bottomrule
\end{tabularx}
\end{table}

%====================================================================
\section{Existing System Architecture}
\label{sec:existing}

\begin{figure*}[!t]
\centering
\fbox{\begin{minipage}{0.95\textwidth}
\centering
\vspace{3em}
\textit{[Figure~1: AHAAR frontend page flow -- /scan -> /detect -> /pantry -> /recipes -> /family -> /cook -> /health. Replace this placeholder with a system-flow screenshot or Mermaid-rendered diagram.]}
\vspace{3em}
\end{minipage}}
\caption{Current AHAAR page flow. Each box corresponds to a route in
the React single-page application.}
\label{fig:flow}
\end{figure*}

AHAAR is implemented as a React~18 single-page application built with
Vite, written in TypeScript, and styled with TailwindCSS. The project
layout is:

\begin{lstlisting}
src/
  App.tsx         # routing
  main.tsx
  pages/          # one file per route
  components/     # reusable UI
  lib/
    data.ts       # demo recipes + ingredients
    quality.ts    # on-device light/sharpness checks
    store.tsx    # client state
  index.css
public/
tailwind.config.ts
vite.config.ts
vercel.json       # SPA rewrites for deep links
\end{lstlisting}

\subsection{Pages and Routes}

\begin{table}[!t]
\caption{Routes of the current AHAAR frontend.}
\label{tab:routes}
\centering
\footnotesize
\begin{tabularx}{\linewidth}{l l X}
\toprule
\textbf{Route} & \textbf{Page} & \textbf{Function} \\
\midrule
\texttt{/}       & Home        & Overview and links to all features \\
\texttt{/scan}   & Fridge scan & Live camera; on-device light and
                                 sharpness checks; capture / upload \\
\texttt{/detect} & Detection   & Bounding boxes with confidence;
                                 remove wrong items \\
\texttt{/pantry} & Pantry      & Editable quantities; ingredient
                                 $\rightarrow$ dish $\rightarrow$ nutrient graph \\
\texttt{/recipes}& Recipes     & Filters; match rings; ``Why this?''
                                 reason bars \\
\texttt{/family} & Family      & Per-person portions by age and
                                 activity, with totals to cook \\
\texttt{/cook}   & Cook-along  & One step at a time; working timer
                                 (keyboard controlled) \\
\texttt{/health} & Health      & Planned CGM-aware portioning
                                 (currently unimplemented) \\
\texttt{/about}  & About       & Pipeline, stack, privacy, team \\
\bottomrule
\end{tabularx}
\end{table}

\subsection{Current Limitations}

\begin{itemize}
  \item Detection boxes are a static demonstration, not live inference.
  \item Only six ingredients and four recipes are supported, all
        hand-coded in \texttt{data.ts}.
  \item Recipe matching uses string equality on ingredient IDs,
        failing on synonyms and partial pantries.
  \item ``Why this?'' bars display hard-coded numbers.
  \item The \texttt{/health} page is a placeholder.
  \item There is no backend, no database, and no model artefact.
\end{itemize}

\S\ref{sec:proposed} addresses each of these limitations.

%====================================================================
\section{Proposed Methodology}
\label{sec:proposed}

\subsection{Fridge-Interior Ingredient Detection}

Existing Indian food datasets~\cite{foodnet2017, indianfood10,
indianfoodnet2023} contain \emph{plated} meals -- not raw ingredients
inside a refrigerator. We therefore construct \textbf{AHAAR-Fridge},
a dataset of 2000+ photographs across 20 ingredient classes commonly
found in Indian kitchens (tomato, paneer, onion, green chilli, coriander,
curry leaves, ginger, garlic, potato, spinach, okra, cauliflower,
carrot, lemon, yoghurt, milk, atta, rice, mustard oil, ghee). Each
image is annotated with YOLO-format bounding boxes.

The detector is a YOLOv11 model pretrained on COCO, fine-tuned first
on IndianFood20~\cite{indianfood10}, and finally fine-tuned on
AHAAR-Fridge. The two-stage transfer protocol follows
\cite{indianfood10, yolov11indian2025}.

The on-device quality gate from \texttt{src/lib/quality.ts} is retained
as a pre-filter -- an element absent from
\cite{redmon2016yolo,foodnet2017,indianfood10,yolov11indian2025}.

\subsection{Semantic Recipe Retrieval}

We scrape approximately 5{,}000 Indian recipes from publicly available
sources (Hebbar's Kitchen, Archana's Kitchen, cookpad). Each recipe
is tokenised and encoded with a multilingual sentence-transformer to
produce a 768-dimensional embedding. The current pantry is encoded by
concatenating detected ingredient names (with their quantities) and
embedding the resulting string. Recipes are ranked by cosine
similarity, following the joint-embedding objective of
\cite{salvador2017im2recipe, marin2021recipe1m}.

\subsection{Nutrition-Weighted Ranking}

Inspired by \cite{yeratoledo2019access}, the final score for recipe
\(r\) and family member \(m\) is
\begin{equation}
\mathrm{score}(r, m) = w_1 \, s_{\text{match}}(r)
                     + w_2 \, s_{\text{diet}}(r, m)
                     - w_3 \, \mathrm{GL}(r)
                     + w_4 \, s_{\text{portion}}(r, m),
\label{eq:rank}
\end{equation}
where \(s_{\text{match}}\) is the pantry--recipe cosine similarity,
\(s_{\text{diet}}\) is a soft dietary-compatibility indicator,
\(\mathrm{GL}(r)\) is the recipe's glycemic load, and
\(s_{\text{portion}}\) is a per-member portion-fit score. Weights
\((w_1,\ldots,w_4)\) are per-member and are updated by the CGM
feedback layer (see \S\ref{sec:cgm}).

\subsection{Set-Based ``Lazy Mode''}

Following Inverse Cooking~\cite{salvador2019inverse}, we add a one-tap
mode that bypasses the pantry edit step: the image is passed to a
set-prediction head that emits an unordered set of ingredients, which
is then used as the pantry input to Eq.~\ref{eq:rank}.

\subsection{Portion Estimation with a Reference Object}

Rather than attempting monocular volume estimation
(explicitly discouraged by~\cite{gao2018bsn}), we require the user to
place a ₹10 coin (26.5~mm diameter) in frame for pantry photographs
that need a true size estimate. The coin provides a pixel-to-millimetre
scale factor, converting detected bounding-box pixel area into a
reasonable gram-weight estimate via per-ingredient density priors.

\subsection{CGM-Aware Personalisation}
\label{sec:cgm}

The \texttt{/health} page accepts CGM exports (Abbott LibreView CSV,
Dexcom Clarity CSV). For each logged AHAAR meal, we extract the
2-hour post-prandial glucose window and compute the area under the
curve (AUC) above baseline. A per-user ridge regression maps meal
macro-nutrient features (carbohydrate load, fibre, fat, protein,
glycemic load) to AUC. Over time the fitted coefficients re-weight
Eq.~\ref{eq:rank} by raising \(w_3\) for users with high glucose
sensitivity. The approach follows the macronutrient-from-CGM
findings of~\cite{huo2019cgm}.

%====================================================================
\section{Backend and Database Design}
\label{sec:backend}

\subsection{Technology Stack}

\begin{itemize}
  \item \textbf{Language:} Java~26 (as of GA, September 2026).
  \item \textbf{Framework:} Spring Boot 3.4 (web, data-jpa, security).
  \item \textbf{Database:} PostgreSQL 16 with Flyway migrations.
  \item \textbf{ML runtime:} Deep Java Library (DJL) 0.30 with the
        PyTorch engine, loading an ONNX-exported YOLOv11 model.
  \item \textbf{Build:} Maven, with a Docker Compose profile for local
        Postgres.
\end{itemize}

\subsection{Database Schema (abridged)}

\begin{lstlisting}[language=SQL]
CREATE TABLE ingredient (
  id        TEXT PRIMARY KEY,
  name      TEXT NOT NULL,
  unit      TEXT NOT NULL,
  kcal_100g NUMERIC,
  carb_100g NUMERIC,
  gi        NUMERIC
);

CREATE TABLE recipe (
  id       TEXT PRIMARY KEY,
  name     TEXT,
  minutes  INT,
  tags     TEXT[],
  embedding VECTOR(768)
);

CREATE TABLE recipe_ingredient (
  recipe_id TEXT REFERENCES recipe(id),
  ing_id    TEXT REFERENCES ingredient(id),
  grams     NUMERIC,
  PRIMARY KEY (recipe_id, ing_id)
);

CREATE TABLE family_member (
  id SERIAL PRIMARY KEY,
  user_id INT, name TEXT, age INT,
  activity TEXT, conditions TEXT[]
);

CREATE TABLE cgm_reading (
  member_id INT, ts TIMESTAMPTZ,
  glucose_mg_dl NUMERIC
);

CREATE TABLE meal_log (
  id SERIAL PRIMARY KEY, member_id INT,
  recipe_id TEXT, grams NUMERIC, ts TIMESTAMPTZ
);
\end{lstlisting}

\subsection{REST API Surface}

\begin{tabular}{ll}
\toprule
\textbf{Endpoint}             & \textbf{Purpose} \\
\midrule
\texttt{POST /api/detect}     & Image upload $\to$ bounding boxes \\
\texttt{GET/PUT /api/pantry}  & Pantry CRUD \\
\texttt{GET /api/recipes}     & Ranked recipes for current pantry \\
\texttt{GET/POST /api/family} & Family members, profiles \\
\texttt{POST /api/cgm/upload} & CGM CSV ingestion \\
\texttt{GET /api/health/me}   & Fitted per-user response model \\
\bottomrule
\end{tabular}

%====================================================================
\section{Dataset Construction and Training Protocol}
\label{sec:dataset}

\subsection{AHAAR-Fridge Dataset}

Images are captured by all three authors with phone cameras under
realistic household conditions: tube-light, LED, warm bulb, and door
light. Each class has 100 training images and 25 validation images.
Images are labelled using the open-source \textit{LabelImg} tool in
YOLO format. The on-device quality gate from \texttt{quality.ts}
rejects blurred or under-exposed captures before they enter the
dataset, reducing label noise.

\subsection{Training Pipeline}

Training is done once in Python/Ultralytics for convenience and the
weights are then exported to ONNX. The ONNX model is served through
DJL at inference time in the Java~26 backend. For transparency, we
also include a DJL-only training loop (ResNet-18 classification head)
to demonstrate end-to-end Java training per the project's learning
goal.

%====================================================================
\section{Evaluation Plan}
\label{sec:eval}

\begin{itemize}
  \item \textbf{Detection:} mAP@50, mAP@50:95, per-class precision and
        recall on AHAAR-Fridge val split.
  \item \textbf{Retrieval:} Recall@1, Recall@5, median rank on a
        held-out 500-recipe query set.
  \item \textbf{Portion:} MAPE in grams on a 100-image coin-reference
        validation set.
  \item \textbf{CGM re-ranker:} RMSE of AUC prediction per user;
        A/B comparison of recipe acceptance rate before and after
        enabling the re-ranker (small user study, $n = 10$).
  \item \textbf{System:} End-to-end latency on a Pixel-class phone and
        on a laptop, broken down per stage.
\end{itemize}

%====================================================================
\section{Discussion}
\label{sec:discussion}

AHAAR deliberately refuses two tempting shortcuts: (i) relying on an
off-the-shelf Indian food detector trained on plated meals, and
(ii) performing monocular volume estimation without a reference
object. The first fails because fridges are a different visual
domain; the second fails because, as \cite{gao2018bsn} documents,
single-image volume is unreliable without depth or scale. The
reference-coin design is honest about this limit.

The CGM-aware re-ranker is the project's most differentiated
contribution. It is the point at which AHAAR stops being yet another
recipe app and becomes a personalised health tool for the Indian
household. No current product combines Indian cuisine recipe retrieval
with CGM-driven glycemic personalisation.

\textbf{Threats to validity.} Our CGM dataset will be small
($<10$ users). The glucose-response model will therefore be
illustrative, not clinical. We will state this limitation explicitly
and target a larger study as future work.

%====================================================================
\section{Conclusion and Future Work}
\label{sec:conclusion}

We have presented AHAAR, a kitchen assistant for Indian households
that spans fridge scanning, ingredient detection, pantry management,
recipe recommendation, per-member portioning, step-by-step cooking,
and CGM-aware health personalisation. The frontend is complete; this
report proposes the backend, ML, database, and dataset work that will
turn the demonstration into a working system. Every design choice is
anchored in a peer-reviewed IEEE paper, and the resulting architecture
is both pedagogically useful (as a Java-26 ML learning project) and
scientifically novel (as a CGM-aware Indian-cuisine assistant).

Future work includes: (i) scaling AHAAR-Fridge to 50 classes with
semi-supervised labelling, (ii) extending the CGM loop to Type-2
diabetic cohorts under IRB-approved supervision, and (iii) a mobile
Android port using TensorFlow Lite.

%====================================================================
\section*{Acknowledgment}
The authors thank Ms.\ Sharlene Rebeiro for her guidance throughout
the project, and the open-source maintainers of Ultralytics YOLO,
Deep Java Library, Spring Boot, and PostgreSQL.

%====================================================================
\bibliographystyle{IEEEtran}
\bibliography{references}

\end{document}
